% Mục 2.3 (người phụ trách: Khang).
\section{Lưu trữ và truy vấn dữ liệu với lakehouse}
\label{sec:2-lakehouse}

Phân tích lịch sử đòi hỏi dữ liệu phải được lưu lâu dài và truy vấn được trên khoảng thời gian dài, điều mà kho gần thời gian thực không đáp ứng được do chỉ giữ dữ liệu 24--48 giờ. Giải pháp lakehouse chia việc này cho bốn thành phần: HDFS lưu các tệp dữ liệu, Apache Iceberg tổ chức các tệp đó thành bảng, Hive Metastore quản lý danh mục để biết mỗi bảng được lưu ở đâu, và Trino là engine truy vấn trên các bảng này.

\subsection{Từ data lake đến lakehouse}
\label{sec:2-lake-to-lakehouse}

Một \emph{data lake} lưu dữ liệu dưới dạng tệp (thường là Parquet \cite{parquet_docs}) trên hệ thống tệp phân tán. Phương án này tốn chi phí thấp và dễ mở rộng, nhưng thiếu các tính năng vốn có của cơ sở dữ liệu quan hệ, như giao dịch, ghi đè nguyên tử, thay đổi schema và truy xuất trạng thái cũ của bảng. Với xử lý micro-batch, vì cứ mỗi 30 giây lại có một lô mới ghi thêm tệp vào bảng, những hạn chế này gây ra hai vấn đề cụ thể. Thứ nhất, nếu một truy vấn đọc bảng đúng lúc lô đang được ghi, truy vấn có thể thấy dữ liệu của lô đó chỉ mới được ghi một phần. Thứ hai, khi một lô phải ghi lại sau sự cố, dữ liệu của lô đó có thể bị ghi trùng.

\emph{Lakehouse} giải quyết điều đó bằng cách đặt một \emph{định dạng bảng} (table format) lên trên tệp Parquet. Định dạng bảng là một lớp siêu dữ liệu cho biết \emph{tập tệp nào hiện là bảng}. Hiện có nhiều định dạng bảng phổ biến như Apache Iceberg \cite{iceberg_spec}, Delta Lake và Apache Hudi, trong đó đề tài sử dụng Apache Iceberg.

\subsection{Apache Iceberg}
\label{sec:2-iceberg}

\paragraph{Cấu trúc siêu dữ liệu.} Iceberg không liệt kê thư mục để biết bảng có tệp nào. Thay vào đó, mỗi bảng là một cây siêu dữ liệu: tệp \emph{metadata} chứa schema, cách phân vùng và danh sách các \emph{snapshot}; mỗi snapshot trỏ tới một \emph{manifest list}; manifest list trỏ tới các \emph{manifest}; mỗi manifest liệt kê các tệp dữ liệu cùng thống kê cột (giá trị nhỏ nhất, lớn nhất, số dòng). Mỗi lần ghi (commit) tạo một snapshot mới và hoán đổi con trỏ tới tệp metadata theo kiểu nguyên tử.

\begin{figure}[H]
\centering
\begin{tikzpicture}[
  node distance=0.5cm and 0.55cm,
  box/.style={rectangle, rounded corners, draw, minimum width=2.3cm, minimum height=0.9cm, align=center, font=\footnotesize, fill=gray!10},
  arr/.style={-{Stealth}, thick}]
\node[box, fill=blue!10] (cat) {Hive Metastore\\con trỏ tới metadata};
\node[box, right=of cat] (meta) {Tệp metadata\\schema, phân vùng,\\danh sách snapshot};
\node[box, right=of meta] (ml) {Manifest list\\của snapshot};
\node[box, below=of ml] (mf) {Manifest\\tệp dữ liệu + thống kê};
\node[box, fill=green!10, left=of mf] (df) {Tệp Parquet\\trên HDFS};
\draw[arr] (cat)--(meta); \draw[arr] (meta)--(ml); \draw[arr] (ml)--(mf); \draw[arr] (mf)--(df);
\end{tikzpicture}
\caption{Cây siêu dữ liệu của một bảng Iceberg}
\label{fig:2-iceberg-tree}
\end{figure}

Cấu trúc này mang lại các khả năng mà thư mục Parquet thuần không có, và mỗi khả năng gắn với một yêu cầu của đề tài:

\begin{itemize}
  \item \textbf{Ghi nguyên tử và cô lập snapshot.} Truy vấn đọc luôn thấy một snapshot hoàn chỉnh, đồng thời lô streaming đang ghi chưa xong không bao giờ lộ ra ngoài. Điều này cần cho các yêu cầu nghiệp vụ BR3 và BR4 (xem Mục~\ref{sec:3-br}).
  \item \textbf{Phân vùng ẩn (hidden partitioning).} Bảng được phân vùng theo một \emph{phép biến đổi} của cột, ví dụ \code{days(measured\_at)}. Người viết truy vấn chỉ lọc theo \code{measured\_at}, Iceberg tự bỏ qua các phân vùng không liên quan. Đây là cơ sở của yêu cầu chức năng FR6.2 (xem Bảng~\ref{tab:3-fr}).
  \item \textbf{Tiến hóa schema.} Thêm, đổi tên hoặc bỏ cột chỉ sửa siêu dữ liệu, không viết lại dữ liệu cũ. Cột được định danh bằng số chứ không bằng tên.
  \item \textbf{Truy vấn theo thời điểm hoặc phiên bản (time travel).} Iceberg giữ nguyên các snapshot cũ nên có thể đọc lại bảng đúng như trạng thái của nó trước đây. Nhờ vậy một báo cáo cũ có thể được tái tạo chính xác, phục vụ yêu cầu nghiệp vụ BR6 (xem Mục~\ref{sec:3-br}).
  \item \textbf{Cập nhật theo dòng.} Phiên bản bảng 2 hỗ trợ xóa và cập nhật theo dòng, nên lệnh \code{MERGE INTO} dùng được cho ghi idempotent (Mục~\ref{sec:2-idempotent}).
\end{itemize}

\paragraph{Vấn đề tệp nhỏ và compaction.} Mỗi micro-batch là một commit, mỗi commit thường sinh ra ít nhất một tệp mới trong mỗi phân vùng nó chạm tới. Với trigger 30 giây, một ngày có khoảng 2\,880 commit; dữ liệu của cả ngày bị chia thành hàng nghìn tệp rất nhỏ. Truy vấn phải mở và đọc siêu dữ liệu của từng tệp nên chậm đi. Iceberg cung cấp thủ tục \code{rewrite\_data\_files} để gộp các tệp nhỏ thành tệp lớn (chiến lược \emph{binpack}) \cite{iceberg_procedures}. Việc gộp tạo một snapshot mới nên không làm gián đoạn các truy vấn đang đọc snapshot cũ. 

\subsection{HDFS làm kho lưu trữ}
\label{sec:2-hdfs}

HDFS (Hadoop Distributed File System) chia tệp thành các khối lớn (mặc định 128 MB), sao chép mỗi khối trên nhiều DataNode và để NameNode giữ danh mục tệp \cite{hdfs_design}. HDFS được tối ưu cho tệp lớn được ghi một lần và đọc tuần tự, nên \emph{không hợp} với hàng nghìn tệp nhỏ: mỗi tệp, dù nhỏ đến đâu, chiếm một mục trong bộ nhớ của NameNode. Đây là một lý do nữa để compaction là bắt buộc chứ không chỉ là tối ưu. Trong đề tài, một NameNode và một DataNode chạy trong Docker; hệ số sao chép là 1 vì chỉ có một nút.

\subsection{Hive Metastore làm catalog}
\label{sec:2-hms}

Iceberg cần một \emph{catalog} để trả lời câu hỏi ``bảng \code{gold.field\_daily} hiện có tệp metadata nào''. Hive Metastore (HMS) là một dịch vụ lưu danh mục bảng và con trỏ tới metadata hiện tại \cite{hive_metastore}. Điểm quan trọng là cả Spark (ghi) và Trino (đọc) cùng đăng ký với \emph{một} metastore, nên cả hai thấy cùng một bảng với cùng một snapshot. Không có catalog dùng chung, hai engine sẽ thấy hai thế giới khác nhau.

\subsection{Trino làm engine truy vấn phân tích}
\label{sec:2-trino}

Trino là engine truy vấn SQL phân tán, không tự lưu dữ liệu mà đọc qua các \emph{connector} \cite{trino_docs}. Connector Iceberg đọc siêu dữ liệu Iceberg qua catalog và đọc tệp Parquet từ HDFS. Để truy vấn dữ liệu phân tích, đề tài dùng Trino thay vì để API gọi trực tiếp Spark với ba lý do sau:

\begin{itemize}
  \item Trino bắt đầu thực thi truy vấn chỉ sau vài mili-giây đến vài giây, trong khi một job Spark phải khởi tạo ngữ cảnh trước khi chạy. Độ trễ khởi động này ảnh hưởng trực tiếp đến thời gian chờ của người dùng trên bảng điều khiển tương tác.
  \item Trino có sẵn máy khách Python (\code{trino.dbapi}) với tham số hóa truy vấn, phù hợp với lớp FastAPI.
  \item Trino đọc được bảng Iceberg \emph{do Spark ghi} mà không cần sao chép dữ liệu, nên không có hệ thống thứ hai phải đồng bộ.
\end{itemize}

Đoạn mã~\ref{lst:2-trino-catalog} là toàn bộ cấu hình catalog của Trino trong đề tài (đã lược chú thích).

\begin{lstlisting}[caption={Catalog Iceberg của Trino}, label={lst:2-trino-catalog}]
connector.name=iceberg
iceberg.catalog.type=hive_metastore
hive.metastore.uri=thrift://hive-metastore:9083
hive.config.resources=/etc/hadoop/conf/core-site.xml,/etc/hadoop/conf/hdfs-site.xml
iceberg.file-format=PARQUET
iceberg.compression-codec=SNAPPY
\end{lstlisting}

Trino mở rộng SQL để dùng các khả năng của Iceberg. Đoạn mã~\ref{lst:2-trino-timetravel} liệt kê các snapshot của bảng bronze bằng bảng siêu dữ liệu \texttt{\$snapshots}, rồi đọc bảng ở trạng thái của một snapshot hoặc một thời điểm.

\begin{lstlisting}[language=SQL, caption={Truy vấn time travel trong Trino (cú pháp của connector Iceberg)}, label={lst:2-trino-timetravel}]
-- cac snapshot cua bang va thoi diem tao
SELECT snapshot_id, committed_at, operation
FROM lakehouse.bronze."sensor_readings$snapshots"
ORDER BY committed_at DESC;

-- doc bang o trang thai cua mot snapshot
SELECT count(*) FROM lakehouse.bronze.sensor_readings
FOR VERSION AS OF 1234567890123456789;

-- doc bang tai mot thoi diem
SELECT count(*) FROM lakehouse.bronze.sensor_readings
FOR TIMESTAMP AS OF TIMESTAMP '2026-11-01 00:00:00 UTC';
\end{lstlisting}

\paragraph{Giới hạn.} Trino không thay thế được kho phục vụ gần thời gian thực như Cassandra. Mỗi truy vấn phải được lập kế hoạch rồi đọc siêu dữ liệu Iceberg và các tệp Parquet nên độ trễ thường từ vài chục đến vài trăm mili-giây (45--120~ms ở p50 trong benchmark \#2), trong khi tra cứu theo khóa trong Cassandra chỉ mất vài mili-giây. Vì vậy hệ thống dùng hai đường phục vụ riêng và không trộn hai nguồn trong cùng một biểu đồ; mức chênh lệch này được đo ở Mục~\ref{sec:5-bench-lookups}.
